\clearpage
\section{Validation}

The repository contains a validation-oriented test suite:

\begin{itemize}[leftmargin=*]
  \item \texttt{FSStatTest}: unit coverage for the shared model, parsing, and
  the three engines on a small fixture;
  \item \texttt{FSStatRobustnessCorrectnessTest}: edge cases and comparative
  validation on identical datasets;
  \item \texttt{FSStatCLITest}: CLI helper and parsing checks.
\end{itemize}

The suite verifies that all three implementations return the same logical
report on:

\begin{itemize}[leftmargin=*]
  \item empty directories;
  \item flat directories with boundary file sizes;
  \item nested hierarchies;
  \item a wider and deeper deterministic hierarchy;
  \item invalid paths;
  \item regular files used as input instead of directories;
  \item unreadable directories, where the environment supports permission
  changes;
  \item symbolic-link cycles, where the environment supports them;
  \item large sparse files;
  \item concurrent file deletions while scanning;
  \item special files and non-regular nodes, where the environment supports
  them.
\end{itemize}

The final verification run executed \texttt{26} tests and all of them passed.

\subsection{Cross-Implementation Oracle}

The main correctness strategy is comparative. For deterministic fixtures, the
Virtual Threads result is compared with the Reactive and Event-Loop results.
The comparison checks the total file count and every histogram band. This is
useful because the implementations have different control structures: if all
three agree on a non-trivial tree, it is unlikely that a scheduling-specific
mistake is hidden in one version.

The tests also validate the shared model directly: band boundaries, duration
formatting, size-unit parsing, and formatted labels are checked independently
from the filesystem traversal.

\subsection{Benchmark Support}

The repository also includes \texttt{FSStatBenchmark}, a small local utility
that runs the three implementations repeatedly and prints timing summaries. It
is useful for manual comparisons between paradigms, but the resulting numbers
are necessarily environment-dependent, so this report does not claim a
universal ordering between the engines.

\subsection{Observed Limitations}

Some checks depend on the host environment:

\begin{itemize}[leftmargin=*]
  \item permission changes may not be applied by the local filesystem;
  \item symbolic links may be unavailable or restricted;
  \item special files such as \texttt{/dev/null} are Unix-specific.
\end{itemize}

These limitations are documented in the validation notes and do not affect the
core functional behaviour of the project.

\subsection{CI-Relevant Stabilization}

Cancellation tests are written to avoid depending on machine speed. In
particular, the Event-Loop cancellation test cancels immediately after the job
handle is returned, instead of waiting for an arbitrary delay and hoping the
scan is still running. This makes the test stable on both local machines and
GitHub Actions.
